\documentclass[UTF8,a4paper,zihao=-4,fontset=none]{ctexart}
\usepackage[margin=24mm,headheight=16pt]{geometry}
\usepackage{fontspec}
\IfFontExistsTF{SimSun}{\setCJKmainfont{SimSun}[BoldFont=SimHei,ItalicFont=KaiTi]\setCJKsansfont{SimHei}\setCJKmonofont{FangSong}}{\setCJKmainfont{FandolSong-Regular}[BoldFont=FandolSong-Bold,ItalicFont=FandolKai-Regular]\setCJKsansfont{FandolHei-Regular}\setCJKmonofont{FandolFang-Regular}}
\IfFontExistsTF{Times New Roman}{\setmainfont{Times New Roman}}{\setmainfont{TeX Gyre Termes}}
\IfFontExistsTF{Consolas}{\setmonofont[Scale=.88]{Consolas}}{\setmonofont{Latin Modern Mono}}
\usepackage{amsmath,amssymb,booktabs,tabularx,array,xcolor,enumitem,fancyhdr}
\definecolor{navy}{HTML}{17324D}\definecolor{teal}{HTML}{087F8C}\definecolor{muted}{HTML}{586579}
\usepackage[colorlinks=true,linkcolor=navy,urlcolor=teal]{hyperref}
\usepackage{xurl}
\hypersetup{pdftitle={CrowdLocateLab：工程复核与开源智能体研究补充报告},pdfauthor={徐楷沅（学号：2352163）}}
\ctexset{section={format=\Large\bfseries\color{navy}},subsection={format=\large\bfseries\color{navy}}}
\setlength{\parskip}{5pt}\linespread{1.13}\setlength{\emergencystretch}{3em}
\setlist{nosep,leftmargin=2em,topsep=5pt}
\newcolumntype{Y}{>{\raggedright\arraybackslash}X}
\pagestyle{fancy}\fancyhf{}\fancyhead[L]{\small\sffamily\color{muted}CrowdLocateLab / 研究与复核}
\fancyhead[R]{\small\color{muted}2026-10-16}\fancyfoot[C]{\small\thepage}\renewcommand{\headrulewidth}{.3pt}
\newcommand{\code}[1]{{\small\nolinkurl{#1}}}
\begin{document}
\thispagestyle{empty}
{\sffamily\large\color{teal}RESEARCH \& REVIEW\hfill 课程项目补充报告\par}
\vspace{5mm}
{\Huge\bfseries\color{navy}CrowdLocateLab\par}
\vspace{4mm}
{\LARGE\bfseries 工程复核与开源智能体研究\par}
\vspace{4mm}
{\large\color{muted}异步测距定位的全局歧义：构造、验证与迭代\par}
\vspace{2mm}
{\normalsize 姓名：徐楷沅\quad 学号：2352163\par}
\vspace{3mm}\noindent\textcolor{teal}{\rule{\textwidth}{1pt}}

\begingroup\small\setlength{\parskip}{3pt}\linespread{1.04}\selectfont
\section*{本版交付结论}
项目已实际接入 Creative-Intelligence 的 Mathematical Concept Agent（Stage 2），由其原始控制器调用模型，在同一会话内完成两个真实研究单元。第一轮提出明确的定位歧义问题、构造反例并编写验证程序；第二轮运行对抗复核、保留失败记录并修复程序。原始控制器与 MIT 许可保留，未把历史自研智能体的运行冒充参考项目 Agent 的成果。

\noindent\begin{tabularx}{\textwidth}{@{}p{28mm}Y@{}}
\toprule
交付项 & 本版提供的证据 \\
\midrule
开源智能体 & 固定上游版本、原始控制器、适配入口、补写 Skills 与真实会话恢复。\\[4pt]
数学研究 & 非共线锚点且局部 Jacobian 满列秩，仍存在相距 6 m 的可行当前位置；两世界最坏误差至少 3 m。\\[4pt]
运行与迭代 & 两轮实际提示、命令、原始事件、状态；v1/v2 程序、第一次失败与后续通过的验证报告。\\[4pt]
可复算成果 & 有理数精确候选检查器、11 个案例、区间 oracle、额外独立复查，以及无需模型的复算入口。\\[4pt]
提交材料 & 本补充 PDF、52 页历史论文、LaTeX 源码、代码与正式证据共同构成课程交付。\\
\bottomrule
\end{tabularx}

\subsection*{成果的准确范围}
本项目研究工业开放题的\textbf{合成数学子问题}。新程序检查给定有限候选，不能自动枚举全部轨迹；\code{no_witness_found} 不证明唯一或安全。锚点、时刻精确及确定性误差硬界是明确假设，尚未由现场数据验证。本报告不主张基础数学结论的文献首创性，也不承诺软件零缺陷。

\subsection*{阅读路线与次数口径}
第 \pageref{sec:upstream} 页起是新增开源 Agent 研究；前面的工程审查保留 2026-10-12 当时的成功、权限错误和修复记录，当前复核见最后一节。历史定位研究共 27 次模型调用、25 个动作、22,052 次定位计算；新增的 2 个研究单元单独计数，单元内工具调用与连接重试不等于 2 次模型后端请求。软件测试与结果重算不增加研究次数。

\vfill\noindent\textcolor{teal}{\rule{\textwidth}{.6pt}}
{\small\color{muted}仓库：\url{https://github.com/orgyuanbao/CrowdLocateLab}\par
本版整理：2026-10-16。发布与云端检查须核对实际提交，不能由 PDF 本身推定。}
\endgroup

\newpage
\section*{历史工程审查：2026-10-12}
本次审查发现了原测试未覆盖的实际缺陷，并在新运行入口、应用检查和证据浏览界面中加入修复与回归测试。可以报告的是\textbf{已检查范围内的结果与剩余限制}；有限测试无法证明程序零缺陷，亦不能保证任意输入的数学正确性或工业部署效果。

本报告对应仓库\href{https://github.com/orgyuanbao/CrowdLocateLab}{orgyuanbao/CrowdLocateLab}，以已发布基线提交\code{270b804}为比较起点。原 52 页论文、冻结估计器、模型日志、数据和测试成绩保持原样。新增运行使用\code{run_research.py}，原入口保留作历史复现。

\subsection*{保持清晰的三条边界}
\begin{enumerate}
\item \textbf{历史科学证据。}两轮共 27 次真实模型调用、25 个研究动作，冻结后共 22,052 次定位计算。这些数字没有因本次检查而增加。
\item \textbf{本次工程验证。}采用离线模型夹具、反例输入、恢复与并发测试，复核既有数据和统计结果；没有新增真实模型研究或独立盲测。
\item \textbf{数值缺陷处置。}不在旧冻结代码中偷偷替换求解器。新应用层独立筛查并拒绝不稳定输出，原估计与校准仅保留作诊断；这不是对任意退化情形的完整求解。
\end{enumerate}

\subsection*{验证证据入口}
完整机器记录保存在\code{verification/hardening/}。其中\code{checks.json}记录软件测试、跳过项、环境、两轮历史审计和统计复核；\code{example_prediction.json}保存新应用入口对项目示例的实际输出。测试总数应以该 JSON 和相应提交的 GitHub Actions 为准，不能把夹具数量当作新科研实验。

首轮工程复核执行 \textbf{184 项测试：183 项通过、1 项因符号链接权限跳过、0 项失败}；比原 129 项增加 55 项。Python 3.12.14、NumPy 2.5.4、SciPy 1.18.1。两轮历史审计均为\code{passed_frozen}，各自零错误、零警告。\textbf{进一步复核发现的缺陷、修复与当时验证见第 \pageref{sec:followup} 页起的记录；上述数字仅属于首轮，不代表本版最终状态。}

\newpage
\section{智能体运行与证据完整性}
\subsection{实际发现及修复}
\noindent\begin{tabularx}{\textwidth}{@{}p{30mm}YY@{}}
\toprule
问题 & 原行为或风险 & 加固后的处理 \\
\midrule
无法冻结时停止 & 规划提示允许停止，工具却拒绝未冻结的 finish，可能无效耗尽预算。 & 允许记录非空停止理由；保持未冻结，不授权留出评测。 \\[4pt]
JSON 数字溢出 & 嵌套参数中的 \code{1e309} 可变成无穷值，后续日志写入失败。 & 响应进入成功日志前检查有限值、重复键和序列化；走受预算限制的错误处理。 \\[4pt]
中断恢复 & 截短事件日志仍可能恢复；单改状态计数可能影响预算或回放位置。 & 先核验链尾、初始化身份、状态与工具快照哈希；不完整请求和不确定工具状态拒绝恢复。 \\[4pt]
回放来源 & 未验证旧事件内容哈希；来源字节及严格提示模式未完全绑定身份。 & 验证完整事件链，以内存快照回放；源文件哈希、严格模式写入运行身份。 \\[4pt]
同目录并发 & 首次启动或恢复竞争可能写入同一日志。 & 首次原子创建目录，使用独占运行锁；崩溃遗留锁要求人工检查。 \\
\bottomrule
\end{tabularx}

\subsection{可恢复不等于任意情况下自动继续}
检查点第 2 版绑定\code{state.json}、工具状态、协议、事件链尾和工具产物目录的字节哈希清单。恢复配置、预算或回放来源改变时拒绝继续。不能自动判断已中断工具的外部副作用是否完成，因此遇到不一致时保留现场，要求核查，而不是重放一个可能已经执行过的操作。

原始两轮目录受新入口保护，且不具有新检查点格式，不能直接冒充新运行恢复。正常退出释放独占锁；进程崩溃的残留锁不能只凭 PID 或时间自动删除。必须先确认原进程已停止、检查证据，再决定后续处理。

\subsection{保证的范围}
哈希是本地一致性检查，不是数字签名。能够协调改写所有记录与哈希的攻击者不在该保证内。真实模型服务、网络、远程模型版本也没有被这些单元测试固定。新入口沿用已公开的开发协议、种子和先验，明确标记为工程运行，不能宣称新盲测。

\newpage
\section{定位应用的数值可靠性}
\subsection{两个可复现反例}
\textbf{坐标原点敏感。}在确定性的 8 锚点无噪声运动夹具中，将所有位置平移至约 $(500000,4000000)$ m 附近，冻结求解器仍报告成功，但可能改变模型选择和位置。局部坐标中的误差约 0.003 m，平移后约 6.959 m。这里比较的是同一物理问题，仅坐标原点改变；它揭示数值终止条件对绝对坐标大小的敏感性，不是新性能基准。

\textbf{弱几何被伪逆截断隐藏。}锚点取 $(10,0),(20,\varepsilon),(30,-\varepsilon)$，目标在原点，使用精确距离。当 $\varepsilon$ 从 $10^{-3}$ 降至 $10^{-4}$ 时，数据更接近共线，但原尺度代理却可能从 100 m 降至约 0.577 m。原因是伪逆截断了弱方向；仅按输出半径判断可能错误接受该结果。

\subsection{应用层增加独立筛查}
\code{application_safety.py}从观测与估计独立重算保留样本的数据 Jacobian，检查参数维度、秩、奇异值、优化成功和活动边界。最小与最大奇异值比值必须满足
\[
\frac{s_{\min}(J)}{s_{\max}(J)}>10^{-5}.
\]
此外将第一个锚点作为新坐标原点，再作一次\textbf{诊断拟合}，比较还原后的点估计、速度和模型选择。位置变化阈值为
\[
\tau_x=\max(0.05\,\mathrm{m},\,10^{-3}L),\qquad
\tau_v=\max(0.001\,\mathrm{m/s},\,\tau_x/\max(\Delta t,1\,\mathrm{s})),
\]
其中 $L$ 为中位距离、最大锚点轴跨度和 1 m 中的最大值。这些阈值是公开的工程规则，并非由留出成绩拟合的最优阈值。非零随机子集初始化因依赖绝对观测哈希而保守拒绝；当前两轮冻结方法该配置均为零。

\subsection{输出语义和成本}
筛查失败时输出\code{status=rejected_numerical_quality}；第二轮应用的\code{acceptance.accepted=false}，即使半径很小或用户未提供半径阈值也一样。原点估计、速度、尺度、区域半径保持不变，不能把诊断拟合结果套入旧校准。

额外拟合次数、工作量与耗时单独记录。原论文速度比较不包含该开销。尚未重新评估新筛查后的接受率、定位误差和条件覆盖；通过筛查仍不能排除全局镜像歧义、模型失配、错误传感器标定或分布偏移。

\newpage
\section{接口、浏览器与验证边界}
\subsection{应用来源与界面容错}
输入先读取为单一字节快照，再由同一快照解析与计算哈希，防止求解期间文件变化导致记录的哈希对应另一份输入。第一轮校准读取也绑定相同快照，且校准水平改为精确匹配，不用近似相等冒充同一水平。

研究浏览器对列表代替对象、空对象、非法嵌套结构、重复键、非有限数字和不完整 JSON 增加容错。无效或缺失证据显示不可用，不冒充零误差。已知文件的解析路径必须位于配置范围内；HTTP 接口保持只读且仅绑定本机。界面存储文本使用文本节点或画布显示，不插入可执行 HTML。它仍是本地研究工具，未认证为生产级多用户服务。

\subsection{如何复查}
在项目根目录安装依赖后执行：
\begin{quote}\ttfamily\small
python -B verification/hardening/run\_checks.py\par
python -B -m unittest discover -s tests -v\par
python audit.py --run runs/revision\_reference --require-frozen
\end{quote}
第一个入口运行软件回归、原始两轮审计及既有统计复核；可用\code{--output}指定不存在的新文件。软件测试包括离线模型夹具和合成反例，不调用真实模型。符号链接测试在缺少 Windows 创建权限时明确跳过，Linux CI 负责执行，不能将跳过写成通过。

本次复核重新计算了既有预测与统计的 123,279 项一致性关系，包括 155 组配对比较与 12 组覆盖总体。它验证已有计算的一致性，没有重新拟合估计器、增加独立样本或创造新模型推理。GitHub 自动检查已覆盖 Windows / Ubuntu 与 Python 3.11 / 3.12，并新增两轮证据审计和统计复核步骤；最终结果需对应本次提交读取。

\subsection{仍需诚实保留的限制}
\begin{itemize}
\item 新主测试静态误差较 WLS 高 29.86\%；90\% 分组区域实际覆盖为 88.17\%；对上一轮总体改善的区间跨零。这些负面结果没有被本次工程加固消除。
\item 研究针对工业开放问题的合成数学子问题，没有真实现场数据，不宣称解决完整工业任务，也不宣称算法首次发明。
\item 数值核心的历史行为保留；新筛查是拒绝风险输出的保护，并非所有数值缺陷均已根治。
\item 有限回归测试、依赖检查与一致性审计都不能推出零漏洞；未来环境变化、未枚举输入和真实模型行为仍需持续验证。
\end{itemize}

\vfill\noindent\textcolor{teal}{\rule{\textwidth}{.6pt}}
{\small\color{muted}提交建议：原论文、此补充报告、全部代码与原始记录一并保留。课堂展示同时说明实际智能体迭代、负面结论、修复证据和适用范围。}
\newpage
\section{进一步复核：证据使用与失败处理}
\label{sec:followup}
\subsection{为什么继续检查}
首轮测试通过后，针对“恢复以后实际读到什么”和“工具只完成一部分时如何处理”补充故障注入。发现两个可复现的运行缺陷，以及一项测试统计缺陷。修复位于新增工程层；原始研究控制器、数值核心、两轮运行日志与 52 页论文均保持原样。

\subsection{实验记录必须在使用时核验}
首轮检查点只绑定状态、工具状态与协议文件，没有覆盖工具状态所引用的实验文件。改动实验 JSON 后，恢复仍可能通过；旧诊断入口还可能返回改动后的误差，同时保留原摘要。这会使后续模型判断依赖与记录不符的数据。

第 2 版检查点增加工具目录的原始字节哈希清单，新增、缺失或改动产物均拒绝恢复。工具调用前核验所引用证据的既有规范化 JSON 摘要；诊断读取一次字节快照，在同一快照上校验和计算，不在校验后重新打开文件。证据路径必须解析在本次运行范围内。

\subsection{工具部分完成后失败必须停止}
多方法实验可能在第一种方法写完文件后，第二种方法计算失败。原处理将异常反馈给模型继续规划，但内存计数、磁盘状态和已写文件可能已经分离，恢复以后会得到另一种状态。

新增包装层比较调用前后的内存状态、持久状态与产物。异常伴随变更时写入不确定失败标记，控制器保留\code{pending_tool}并设为\code{tool_failure_uncertain}，停止后续模型调用并拒绝自动恢复。标记写入失败时仍保留控制器待定状态。验证参数时、尚未产生副作用的普通错误仍可作为观察反馈。已有不完整运行需要人工检查，不能删除标记或重写哈希以冒充一致。

无法建立调用前状态快照时也保守停止；命令行对该不确定失败返回退出码 2，避免调用者误判为正常完成。

\subsection{检查点兼容性与信任边界}
运行配置中的保护版本更新为 3，检查点格式更新为 2；缺失或旧格式检查点明确拒绝恢复，不静默升级旧记录。新工程运行必须使用新目录。回归覆盖文件篡改、缺失、额外产物、仅格式变化、诊断快照、路径越界、内存或磁盘部分变更、存储故障与正常无副作用错误。

哈希校验用于发现不一致，不能防止具有同等文件写权限的攻击者同时伪造所有证据与摘要，也不是面对任意并发恶意写入的事务文件系统保证。科研结果的可信范围仍需结合完整日志与原始数据理解。

\newpage
\section{统计勘误、数学公式与最终复查}
\subsection{通过数不能由错误事件数倒推}
\code{unittest}中的同一测试可产生多个子测试错误。原复核脚本以“运行数减错误数减失败数”计算通过数，在环境权限造成多个子测试报错时出现负数。现改为只在父测试实际成功时计数；失败、错误和跳过仍保留框架的事件数，不能与父测试数直接相减。回归包括多子测试失败、多子测试错误、跳过和预期失败情形。

首次重跑还暴露当前环境对默认临时目录的写入限制。最终复查使用工作区内可写临时目录；环境失败记录保留在本次复核工作目录，没有改写成通过，也没有以缩减测试范围解决权限问题。

\subsection{原论文公式的实现勘误}
原论文距离残差部分将尺度简写为 $q_i=\sigma_i$。冻结的组合估计器实际先计算尺度，再施加 $10^{-5}$ 的数值下界。因此在\code{residual=range, scale=reported}配置下，准确表达为
\[
q_i=\max(\sigma_i,10^{-5}).
\]
例如合法输入 $\sigma_i=10^{-8}$ 时，程序使用 $10^{-5}$。对于\code{scale=constant}且没有显式常量的配置，常量取施加下界后的尺度中位数。该说明仅对应组合估计器的尺度实现，不把其他历史基线一概等同。这里更正文献表述，\textbf{没有改变代码、历史拟合值、校准半径或性能成绩}。

\subsection{最终验证与记录}
最终机器复核记录为\code{verification/hardening/final_checks.json}，包含软件回归计数、运行环境、两轮完整历史审计、冻结状态和既有统计一致性检查。本轮未调用真实模型、未生成新的科研预测，也没有以软件夹具冒充研究迭代。

\textbf{本轮 208 项：205 项通过、0 项断言失败、2 项权限错误、1 项跳过，整体仍为未全通过。}新增 24 项回归全部通过。两个 HTTP 集成用例在连接本机服务时被当前网络权限以\code{WinError 10013}拒绝；一项符号链接夹具受系统权限限制。它们没有验证完成，必须在具备权限的环境或对应提交 CI 补验。诊断保存在\code{verification/hardening/final_software_diagnostics.txt}。

两轮完整历史审计均为\code{passed_frozen}且零错误、零警告；冻结检查及 123,279 项统计一致性检查通过。原 52 页论文与冻结研究证据保持原样。

发布状态以相应提交和 GitHub Actions 为准。此补充报告不能单独证明远端已更新，不能用先前提交的绿色状态替代当前提交的检查结果。原论文中的静态退化、覆盖不足、区间跨零以及缺少现场验证等结论继续保留。

\vfill\noindent\textcolor{teal}{\rule{\textwidth}{.6pt}}
{\small\color{muted}历史记录到此结束。以下为 2026-10-15 的真实 Agent 研究和 2026-10-16 的交付复核。}

\newpage
\section{直接采用参考项目的 Agent}
\label{sec:upstream}
\subsection{来源、接入与角色分工}
参考项目为 \href{https://github.com/DeepMathLLM/Creative-Intelligence}{DeepMathLLM/Creative-Intelligence}，固定于提交\code{b7d17249f94949351faefd7ae28a020f0623db93}。本项目逐字节保留其 Stage 2 的\code{agent.py}、配置、说明、模式、测试及 MIT 许可，共 7 份来源文件。其原始控制器实际执行\code{command_run}、\code{execute_turn}与会话恢复；适配器只提供研究工作目录、配置、输入摘要及过程证据。

该公开快照未包含 README 所提到的构造、验证两份 Skill，因此本项目补写\code{construct-mathematical-concept}与\code{verify-mathematical-concept}，逐份注明归属。它们是研究指引，不是形式化证明器。Stage 2 使用 Codex CLI，不等同于该项目第一阶段的 Moonshine 运行时。

模型可以在独立研究目录内写数学推导与 Python 程序，执行计算并根据反馈修订。历史自研控制器的受限工具研究与本次入口分开：前者保留旧两轮定位成绩，后者允许形成新的验证程序。题面、参考材料与 Skills 的输入摘要在调用前后检查。

\subsection{两次真实研究推进}
\noindent\begin{tabularx}{\textwidth}{@{}p{25mm}YY@{}}
\toprule
记录 & 第一次推进 & 第二次推进 \\
\midrule
日期 & 2026-10-15 & 2026-10-15 \\[3pt]
实际任务 & 构造、推导、编写检查器并运行案例 & 同会话恢复；尝试反驳、运行 oracle、修复工程问题 \\[3pt]
接受状态 & turn=1，paused & turn=2，complete \\[3pt]
主要产物 & \code{research_v1.md}、\code{diagnostic_v1.py}、11 个案例及结果 & \code{diagnostic_v2.py}、\code{verification_v2.md}、失败与通过报告 \\
\bottomrule
\end{tabularx}

两次\code{accepted_structured_research_turn}均为 true，保存状态中的 session ID 相同，第二次实际命令包含\texttt{exec resume}。原始 JSONL 事件保留了读写文件、执行程序和连接重试；没有把模拟进程测试计为真实研究。\code{complete}只对应题面限定的见证研究，不认证数学正确性或工业题已解。

\subsection{证据与运行边界}
原始记录位于\code{studies/creative_intelligence/mirror_ambiguity/evidence/}。每轮保留实际提示、命令、状态前后快照、原始事件、错误输出与摘要；第一轮科学产物另存快照。SHA256 用于核查复制与变化，不是不可伪造的第三方签名。

每次\code{advance}仅推进一个研究单元，\code{--once}不构成令牌、费用或时间上限。运行采用独立可写工作区，禁用研究沙箱网络和递归子 Agent；模型服务的连接仍由 CLI 完成。停止机制不能保证递归终止所有后代进程，故本次未启动后台或递归任务。

\newpage
\section{问题定义与精确判据}
\subsection{固定观测与存在共同观测}
令当前位置与速度为 $\theta=(x,v)\in\mathbb R^4$，锚点为 $a_i$，时刻为 $t_i\le0$，并记
\[
d_i(\theta)=\|x+t_i v-a_i\|,\qquad
\mathcal F(y,\epsilon)=\{\theta:|y_i-d_i(\theta)|\le\epsilon_i,\ \forall i\}.
\]
新研究假设锚点、时刻精确，$y_i\ge0$，$\epsilon_i\ge0$ 是\textbf{确定性误差硬界}，不是历史实验的标准差 $\sigma_i$。给定实际 $y$ 时，须找到 $\theta_1,\theta_2\in\mathcal F$ 且 $x_1\ne x_2$ 才形成当前位置歧义见证。

另一问题是：对指定的两个候选，是否存在某个共同观测向量。这等价于逐行
\[
|d_i(\theta_1)-d_i(\theta_2)|\le2\epsilon_i.
\]
逐行中点给出共同观测，但它未必等于实际输入 $y$。例如两距离为 1、3 且 $\epsilon=1$，共同观测 2 可行；实际观测 1 对第二个候选不可行。

\subsection{避免浮点临界误判}
程序对输入有理数精确计算平方距离 $q_i=d_i(\theta)^2$。固定观测判据为
\[
\max(0,y_i-\epsilon_i)^2\le q_i\le(y_i+\epsilon_i)^2.
\]
对两个候选的平方距离 $q,r\ge0$，共同观测模式设 $S=q+r-4\epsilon^2$，其等价判据为
\[
S\le0\quad\text{或}\quad (S>0\ \text{且}\ S^2\le4qr).
\]
第二次平方只在两侧非负时使用，不能删除 $S\le0$ 分支。输出同时提供位置分离平方 $D^2$ 与误差下界平方 $D^2/4$；显示的有理下界向下取整，不能向上夸大。

\subsection{近共线与异步反射}
以经过 $c$、单位法向为 $n$ 的直线反射整条轨迹，记
\[
s_i=n\cdot(x+t_i v-c),\qquad h_i=n\cdot(a_i-c).
\]
两反射候选的距离差在分母非零时满足
\[
|d_i-d_i'|=\frac{4|h_i s_i|}{d_i+d_i'}\le2\min(|h_i|,|s_i|).
\]
双零距离时差为零，不能代入 $0/0$。逐行 $h_i s_i=0$ 等价于精确反射测距相同。近共线给出的界是保守充分条件；异步过线时 $s_i=0$，离轴锚点也可能无法消歧。速度只作线性反射，位置包含仿射平移。

\newpage
\section{非共线、满秩仍有全局歧义}
\subsection{Agent 构造的精确反例}
取 $\theta_1=(0,3,1,1)$、$\theta_2=(0,-3,1,-1)$，坐标依次为 $(x_1,x_2,v_1,v_2)$。每条观测的误差硬界均为零：
\begin{center}
\begin{tabular}{@{}rrrr@{}}
\toprule
$t_i$ & 锚点 $a_i$ & 两条轨迹的距离 & 行号 \\
\midrule
$-4$ & $(-4,0)$ & $1$ & 0 \\
$-3$ & $(-3,1)$ & $1$ & 1 \\
$-2$ & $(-2,0)$ & $1$ & 2 \\
$-1$ & $(-5/2,0)$ & $5/2$ & 3 \\
$0$ & $(-4,0)$ & $5$ & 4 \\
\bottomrule
\end{tabular}
\end{center}
前三锚点的有向二倍面积为 $-2\ne0$，确实非共线。离轴锚点对应 $t=-3$，两条轨迹恰好在该时刻经过反射轴，因而距离仍相同。该例在\code{fixtures_v1.json}中名为\code{noncollinear_crossing_full_rank}。

令 $z_i=x+t_i v-a_i$，距离 Jacobian 的第 $i$ 行是 $[z_i^\top,t_i z_i^\top]/d_i$。所有距离非零，因此行乘 $d_i$ 不改变秩。第一条轨迹取第 0、1、3、4 行得到
\[
B_1=\begin{pmatrix}
0&-1&0&4\\
0&-1&0&3\\
3/2&2&-3/2&-2\\
4&3&0&0
\end{pmatrix},\qquad \det B_1=-6.
\]
镜像轨迹对应子式行列式为 $+6$。两个 Jacobian 都满列秩 4，但当前位置分别为 $(0,3)$ 与 $(0,-3)$，分离距离 $D=6$。所以局部满秩及非共线锚点不足以保证全局唯一。

\subsection{为什么存在 3 m 的不可消除下界}
同一观测产生的任一确定性单点估计 $z$ 都满足
\[
6=\|x_1-x_2\|\le\|z-x_1\|+\|z-x_2\|,
\qquad \max_j\|z-x_j\|\ge3.
\]
因此不存在对这两个可行世界都保证小于 3 m 误差的估计。对相同输入采用同一输出分布的随机估计，取期望同样得到两世界中较大的期望误差至少为 3。这里不是概率覆盖率结论；在 $D=2r$ 时，中点可达到误差 $r$，只有 $D>2r$ 才否定所有世界误差不超过 $r$ 的保证。

此结果是具体验证见证，不主张完整唯一性判定或数学文献首创。缺少候选、漏掉镜像时，程序仍只能报告未找到见证。

\newpage
\section{真实失败、修复与复核}
\subsection{第二轮实际发生的迭代}
第二轮使用同一模型会话，另外实现区间 oracle，直接比较原始允许观测区间，不复用待检平方判据生成预期答案。有限计算支持推导，但不是普遍命题的形式化证明；另写 oracle 也不等于不同模型的确认。

第一次验证退出码为 1，保存为\code{results_v2.json}：验证脚本错误地要求两轨迹的子式均为 $-6$，而镜像应为 $+6$；同轮还出现 Windows 子进程输出解码错误。修正验证脚本与 UTF-8 调用后，保存\code{results_pass_v2.json}；补充零距离和字节边界后，再保存\code{results_final_v2.json}。失败没有被重标为成功，11 份 v1 科学文件的 SHA256 未变。

\noindent\begin{tabularx}{\textwidth}{@{}p{30mm}YY@{}}
\toprule
确认的程序问题 & v1 的安全复现 & v2 的处理 \\
\midrule
重复 JSON 字段 & 后一个值静默替代前一个 & 解析阶段拒绝任意对象中的重复字段 \\[4pt]
覆盖输入或结果 & 输出可以写回输入或旧结果 & 拒绝相同路径与已有文件，采用排他新建 \\[4pt]
无输入规模限制 & 温和超界输入仍被接受 & 限 32768 字节、32 条观测、16 个候选；限制字面长度、有理数大小与嵌套 \\
\bottomrule
\end{tabularx}
数值界在构造\code{Fraction}前检查巨大指数；支持域外输入拒绝属于工程范围收缩，不是数学唯一性判定。排他新建防止覆盖，但不保证磁盘故障时新文件完整写入。

\subsection{已完成的有限检查}
\begin{itemize}
\item 每个版本 517 个实际观测判据案例、514 个共同观测判据案例，与区间 oracle 零不一致；包含阈值两侧、等号、无理根夹逼与零距离。
\item 11 个既有案例及其旋转平移、85 个反射案例、满秩子式、$D/2$ 下界、输出向下界均完成检查。
\item 最终 35 次 CLI 调用均符合预期状态，包含应当拒绝的输入，不能描述为 35 次都退出 0。
\item 项目另一个审阅 Agent 对实际产物复查：第一版 9000 个精确算术案例；第二版 20 项数值边界、9 项 CLI、维度与输出保护等通过。该复查仍不是外部同行评议或形式化证明。
\end{itemize}

\subsection{可移机复算}
\code{reproduce_checks.py}先核验 70 份原始材料的摘要，在临时副本中运行原验证及第二版独立回归，仅重定位历史哈希清单的绝对路径，不改归档原件、不调用模型。默认只需 Python 标准库；需要 NumPy 的第一版外部复查保留原结果，不包含在此默认入口中。

\newpage
\begingroup\small\linespread{1.04}\selectfont
\section{本版交付核对与剩余问题}
\label{sec:delivery}
\subsection{可复现入口}
从项目根目录执行\texttt{python -B}\ \code{studies/creative_intelligence/mirror_ambiguity/reproduce_checks.py}，可复算新增研究。使用\code{--output}必须指定尚不存在的新目录。实际模型研究则通过\code{integrations/creative_intelligence/runner.py}的\code{init}与\code{advance}执行，需要自己的 CLI 登录和额度；再次运行可能产生不同输出，不覆盖正式记录。

历史审计使用\texttt{python audit.py --run}\ \code{runs/live_reference}\ \texttt{--require-frozen}，再将目录替换为\code{runs/revision_reference}检查第二轮。\code{verification/check_revision_results.py}检查已保存统计，不重新拟合。全部命令、数值环境和真实输出以本版\code{studies/creative_intelligence/mirror_ambiguity/verification/}及对应提交 CI 为准。

\subsection{2026-10-16 的完整检查}
首次完整本地回归执行 \textbf{273 项：271 项通过、2 项跳过、0 项失败或错误}，包含原工程 233 项、上游控制器 22 项与适配层 18 项。两项跳过因 Windows 无符号链接权限，须由对应提交 Linux CI 补验。

首轮云端检查中 Windows 两组通过，Linux 发现上游 3 个显示测试固定断言 Windows 的\texttt{PowerShell}，而 Linux 正确显示\texttt{Shell}。现仅在这些测试夹具内提供明确的 Windows 平台上下文，并增加 2 个宿主与显式命令标签检查；原控制器、原测试文件和数学产物未改。更新后测试总数为 275，本地新增两项通过；云端状态以修正后的提交为准，失败记录保留。

两轮历史审计均为\code{passed_frozen}，零错误、零警告；123,279 项统计一致性检查通过。新增研究的 6 组复算、35 次 CLI 回归和第二版独立检查通过，70 份原始材料摘要保持不变；两轮历史归档成功重建为 14 张和 13 张卡片。环境为 Python 3.12.14、NumPy 2.5.4、SciPy 1.18.1，以上收尾检查新增模型调用为零。

原始命令、标准输出与检查汇总保存在\code{verification/final_20261016/}（相对于新研究目录）。PDF 由同一份 LaTeX 源码编译并作逐页检查；云端四组环境必须对应最终上传提交。报告不以进程退出成功替代科研正确性，不以历史绿色状态替代新提交结果。

\subsection{课程要求与成果归属}
智能体源码、数学问题与反例、两轮真实调用、失败与修订、PDF 与复算代码均有对应证据，封面列出提交映射。模型实际编写并运行新程序，作者适配与上游来源分别注明。本补充报告与 52 页历史论文共同提交；学生为徐楷沅（2352163）。

\subsection{尚未解决的研究问题}
\begin{itemize}
\item 如何自动发现候选，或完整判定连续参数空间中的全局唯一性。
\item 如何为真实系统标定误差硬界，并处理锚点位置、时间戳的不确定性。
\item 如何将歧义见证转化为完整的应用拒绝策略并验证拒绝后的精度、成本与覆盖。
\item 如何获得现场数据与完整工业验收标准。原统计研究的静态退化和覆盖不足仍须保留。
\end{itemize}

\vfill\noindent\textcolor{teal}{\rule{\textwidth}{.6pt}}
{\small\color{muted}本项目交付的是有真实运行、有失败迭代、有数学与程序证据的课程研究。检查通过说明已检查范围内未发现阻断问题，不等于任意输入、未来环境或完整工业问题均已正确解决。}
\endgroup
\end{document}
